% Scratch file for Claude's computations, derivations, and intermediate work.
% Review results here before incorporating into the main writeup.

% ============================================================
% Professional rewrite of Section 7.1, lines 11-23
% (Recap: Perturbative String Theory and Higher-Derivative Supergravity)
% ============================================================

\subsection{Recap: Perturbative String Theory and Higher-Derivative Supergravity}
Higher-derivative corrections to supergravity are most efficiently obtained not from Feynman diagrams, but from the low-energy expansion of string amplitudes.
The perturbative expansion of string amplitudes is organized by two parameters: $\alpha'$, governing the derivative expansion of the low-energy effective action, and $g_s$, weighting the genus expansion of the worldsheet. An $n$-point amplitude admits the expansion
\begin{equation}
    \mathcal A_n = \sum_{g=0}^\infty g_s^{2g-2+n}\, \mathcal A_n^{(g)}(\alpha').
\end{equation}

% ============================================================
% Derivation: direct evaluation of the 4-point KK-loop R^4 coefficient
% (the step "shortcut" in Sec 7.3; full version goes to Appendix A__KK_loop_R4)
% ============================================================
% Leading R^4 coefficient = box with KK tower, low-energy limit P(0,0;lambda)=1/6 (simplex vol):
%   C = sum_{(m,n)!=0} int_0^inf dlambda lambda^{-3/2} exp(-lambda M^2),  M^2=|m+n tau|^2/(R^2 tau2^2)
% Proper-time integral via analytic continuation (zeta reg discards power-UV = higher-d R^4):
%   I(s) = sum int dlambda lambda^{s-1} e^{-lambda M^2} = Gamma(s) R^{2s} tau2^s E_s(tau)
%   with E_s = sum_{(m,n)!=0} tau2^s/|m+n tau|^{2s}.   Our case s=-1/2.
% Functional eq (Poisson resummation): pi^{-s}Gamma(s)zeta(2s)E_s = invariant under s->1-s
%   s=-1/2 -> E_{-1/2} = (3 zeta(3)/pi^2) E_{3/2}.   [verified numerically: 0.36538...]
% Constant terms of E_{3/2} = 2 zeta(3) tau2^{3/2} + (2 pi^2/3) tau2^{-1/2} + instantons = f_0.  [verified]
% => C = Gamma(-1/2) R^{-1} tau2^{-1/2} (3 zeta(3)/pi^2) f_0 = -(6 zeta(3)/(pi^{3/2} R tau2^{1/2})) f_0  ~ f_0.
% DIMENSIONAL REMARK: lambda^{3-d/2}; d=9 (11d on T^2) -> lambda^{-3/2} -> E_{3/2}=f_0 (correct IIB weight).
%   d=10 (genuine 12d on T^2) -> lambda^{-2} -> I(-1) -> E_{-1} -> E_2 (wrong weight) + Gamma(-1) pole (log).
%   Hence the "spectator dimension" shortcut works because the result is intrinsically 9d;
%   reinforces that t8t8R^4 is not intrinsically 10d. True 10d diagnostic: eps10 eps10 R^4 at 5-pt.
